\uuid{oZ1n}
\exo7id{2316}
\titre{exo7 2316}
\auteur{matexo1}
\organisation{exo7}
\datecreate{2002-02-01}
\isIndication{false}
\isCorrection{false}
\chapitre{Calcul d'intégrales}
\sousChapitre{Théorie}
\cpge{mpsi-mp2i.integration.c, pcsi.integration.b, ptsi.integration.b, mpsi-mp2i.espaces-prehilbertiens.b, pcsi.espaces-prehilbertiens.b, ptsi.geometrie-espace.b, bcpst1.integration-segment.a, bcpst1.integration-segment.b, mp-mpi.fonctions-vectorielles.c}
\module{Analyse}
\niveau{L1}
\difficulte{}

\contenu{
\question{
Soit $f$ une fonction continue sur l'intervalle
$[a,b]\subset\Rr$, \`a valeurs dans $\left]0,+\infty\right[$. Montrer
que
$$ \int_a^b f(t)\,d t \times \int_a^b \frac 1{f(t)}\,d t
\geq  (b-a)^2.$$
Dans quels cas y a-t-il \'egalit\'e ? (On pourra utiliser
l'in\'egalit\'e de Cauchy-Schwarz.)
}
}
